\documentclass[acmlarge]{acmart}

\usepackage{booktabs}
\usepackage{tabularx}
\usepackage{array}
\usepackage{multirow}
\usepackage{makecell}

\AtBeginDocument{%
  \providecommand\BibTeX{{%
    Bib\TeX}}}



\begin{document}

\title{Arm Cortex-A76 Validation Study via gem5}

\author{Eric Rippey}
\email{erippey3@vt.edu}


\begin{abstract}
This paper investigates the use of gem5 to approximate and validate the CPU
behavior of the Broadcom BCM2712, the system-on-chip used in the Raspberry Pi 5
family. The BCM2712 contains a quad-core Arm Cortex-A76 processor, a class of
processor widely used in modern embedded systems. In this work, a gem5 simulation
of the BCM2712 is constructed to emulate key aspects of the processor and is
evaluated on multiple workloads. The study compares results from both the
simulated system and a physical Raspberry Pi 5 to identify where the simulation
matches hardware behavior and where it diverges.
\end{abstract}




\received{March 23, 2026}

\maketitle

\section{Introduction}
Modern embedded and edge systems increasingly rely on processors that balance
performance, power efficiency, and cost. Among these processors, the Arm
Cortex-A76 has seen adoption in embedded system-on-chips from vendors such as
Broadcom, Qualcomm, Rockchip, and Samsung. A commonly accessible example is the
Broadcom BCM2712 used in the Raspberry Pi 5 and Raspberry Pi 5 Compute Module.
Because platforms in this class operate under tight power and thermal
constraints, understanding their performance characteristics is important for
both system designers and application developers.

On physical hardware, tools such as \texttt{perf} and vendor-specific profilers
such as Arm Streamline Performance Analyzer can expose detailed performance
behavior through hardware counters. These tools are convenient, but physical
hardware has a fixed design and is subject to experimental noise. Simulation
provides a controlled and repeatable environment in which hardware parameters can
be modified for theoretical studies. Architectural simulators such as gem5 are
therefore attractive because they provide access to internal statistics and
enable experiments that would not be possible on fixed hardware. However, their
usefulness depends on how faithfully they reproduce real system behavior.

The desire to faithfully approximate the BCM2712 is the central motivation of
this study. We aim to understand the extent to which gem5 can model the behavior
of the BCM2712's Cortex-A76 CPU subsystem by comparing it to a physical
Raspberry Pi 5. We evaluate whether a gem5 model can reproduce enough of the
target platform's behavior to support meaningful architecture experiments. Thus,
this work is a validation study. If the simulated BCM2712 tracks hardware metrics
closely enough, it can become a credible tool for future exploratory studies. If
it does not, the differences can help identify the limitations of simulating this
platform.

\section{Background}
The gem5 simulator was first released in 2011 and was designed to merge the best
aspects of the M5 and GEMS simulators\cite{Binkert2011-oh}. These strengths
included high configurability, support for multiple instruction set
architectures, flexible memory systems, and multiple memory interconnect models.
The gem5 simulator supports several CPU models. It provides \texttt{AtomicSimple},
a single-instruction-per-cycle model; \texttt{TimingSimple}, a
single-instruction-per-cycle model that simulates memory reference timing;
\texttt{InOrder}, a pipelined in-order execution model; and \texttt{O3}, a
pipelined out-of-order execution model. Applications in gem5 can be run in either
\texttt{system-call emulation} mode, which avoids the need for an operating
system, or \texttt{full-system} mode, which executes both user and kernel-level
instructions. The gem5 platform also includes two memory systems: the
\texttt{Classic} model, which is fast and straightforward to configure, and the
\texttt{Ruby} model, which provides flexible configurability for different cache
coherence strategies.

Arm released a research starter kit for modeling Arm architectures using
gem5\cite{githubGitHubArmeducationarmgem5rsk}. The Arm kit includes guided
documentation and a GitHub repository. In addition to defining many CPU and
memory options provided by gem5, the starter kit provides a High-Performance
In-Order CPU with timings tuned to represent a modern in-order Armv8-A
implementation. This differs from the BCM2712's Cortex-A76, which is an
out-of-order core. The High-Performance In-Order CPU contains four stages: two
fetch stages, a decode stage, and an execute stage. For each stage, parameter
specifications are provided to tune pipeline behavior so that it resembles an
Armv8-A implementation. Similar specifications are provided for branch
prediction, specialized functional units, level 1 instruction and data caches,
the TLB, and L2 cache. The kit does not provide specifications for L3 cache,
which is optional in Cortex-A76 systems\cite{armCompare}. This prior work is
limited to a generic Armv8-A in-order architecture. This study instead implements
an out-of-order Armv8-A model to better approximate the Cortex-A76\cite{armCortexA76Product}
and provides validation results against physical hardware.

There have been previous attempts to recreate Arm architectures in gem5 and
validate them against physical hardware. Butko et al. compared a simulated Arm
Cortex-A9 against an ST-Ericsson Nova A9500 dual-core Arm Cortex-A9 processor
\cite{Butko2012-vg}. Their study compared the simulated two-core CPU to the
physical device across 11 benchmarks with varied execution profiles. Benchmark
runtime differed by 0.47 to 17.94 percent between the real and simulated systems,
with a mean difference of 6.4 percent. The paper also examined dense matrix
operations and noted that gem5 struggled to accurately simulate L2 cache misses,
overreported misses, and used a DDR memory model that did not accurately
represent the tested system. This paper conducts a similar study for the Arm
Cortex-A76 with two key differences. First, it examines additional performance
counters beyond execution time and L2 cache misses. Second, it implements an
out-of-order CPU model to better match the Cortex-A76 pipeline.

In another study, Endo et al. simulated the Arm Cortex-A8 and Cortex-A9 using
in-order and out-of-order pipelines, respectively\cite{Endo2014-ca}. Both
architectures were modeled using the O3 CPU model; for the A8, the dispatch order
was modified and the register renaming stage was disabled. The study collected as
many available parameters as possible from Arm technical manuals, but the authors
were unable to find information about L1 cache buffer sizes, some branch
predictor parameters, TLB sizes, instruction queue sizes, and reorder buffer
sizes. The paper clearly defines the parameters it found and the assumptions it
made. The simulated CPUs were compared against physical Cortex-A8 and Cortex-A9
systems using 10 PARSEC 3.0 benchmarks in gem5's \texttt{full-system} mode. The
authors found runtime differences of 1 to 17 percent for the A9 and 2 to 16
percent for the A8. They attributed most of the error to the general-purpose
nature of gem5 and the need to use default parameters where implementation
specific values were unknown. This study contributes similar work for the Arm
Cortex-A76 and investigates performance counters beyond runtime. While the prior
paper focused on runtime and data movement between VFP/NEON processors and the
Arm pipeline, this study also includes cache statistics and branch behavior.

There have also been attempts to model specific parts of the Arm Cortex-A76. In
particular, Schluter et al. used side channels observed during benchmarks to
analyze hardware prefetcher characteristics\cite{FetchBench}. Through
experimentation, they identified characteristics of 19 different Arm and x86-64
CPUs, including the Arm Cortex-A76. Arm does not provide explicit details about
the Cortex-A76 hardware prefetcher in the Technical Reference Manual
\cite{armCortexA76Product}. The manual describes supported preload instructions
and states that the load-store unit contains a prefetcher responsible for both
L1 and L2, but it does not provide the underlying prefetcher design. To infer
this withheld information, Schluter et al. present FetchBench, a benchmark suite
with distinct memory access patterns. Each test case is characterized on a given
processor and compared against the expected behavior of common prefetching
algorithms.

\section{Methodology}

This study evaluates how closely a gem5 recreation of the Arm Cortex-A76 can
approximate the behavior of the Raspberry Pi 5 CPU. The target reference platform
is the Raspberry Pi 5, which uses the Broadcom BCM2712 SoC containing four Arm
Cortex-A76 cores. The goal of the model is to construct a reproducible
approximation of the CPU, cache hierarchy, and memory subsystem that captures
behavior as close to the physical CPU as possible. Using this approximation, the
study quantifies where the recreation reproduces hardware behavior and where the
gem5 model diverges.

\subsection{Reference Platform}

The reference hardware used in this study was an 8 GB Raspberry Pi 5 based on
the Broadcom BCM2712 SoC. The BCM2712 includes four Arm Cortex-A76 cores and
LPDDR4X memory. This platform represents a modern embedded Arm system with an
out-of-order Cortex-A profile processor. The Cortex-A76 contains a 13-stage
pipeline, two simple ALUs, one multi-cycle ALU, one branch unit, two 128-bit
ASIMD/FP pipelines, and two load/store units\cite{CortexOptimizationGuide}. By
comparison, the gem5 O3 CPU is based on the Alpha 21264, which has a seven-stage
out-of-order execution model\cite{584667}.

For hardware measurements, the Raspberry Pi 5 was configured to reduce
experimental noise where possible. The CPU frequency was fixed at 2400 MHz so
that the hardware could be compared against a simulated processor running at the
same nominal frequency. Background services were minimized where possible.

\subsection{Simulation Scope}

The gem5 model focuses on the CPU system rather than the full Raspberry Pi 5
system-on-chip. Specifically, the simulation attempts to approximate the
Cortex-A76 cores, private L1 instruction and data caches, shared L2 cache, L3
cache, system memory, and TLB behavior.

\subsection{Model Construction and Parameter Selection}

The simulated architecture was constructed using the \texttt{ArmO3CPU} model.
Progress toward an accurate model of the Arm Cortex-A76 proceeded in stages. The
first iteration of the \texttt{ArmO3CPU} was based largely on gem5 default
values. Iteration 1 consisted of four out-of-order CPUs, three cache levels sized
to match the Cortex-A76-based Raspberry Pi 5 platform, and generic values for all
other parameters. This iteration served as a baseline for all subsequent changes.

The second iteration of the simulated CPU was gradually tuned using publicly
available documentation and justified approximations. Parameters were selected
using the following priority order:

\begin{enumerate}
  \item Public Arm Cortex-A76 documentation and technical reference material \cite{armCortexA76Product,CortexOptimizationGuide}
  \item Public Arm A-profile architecture documentation and technical reference material \cite{arm-a-documentation}
  \item Raspberry Pi 5 and BCM2712 documentation where available
  \item Prior gem5 and Arm simulation studies
  \item Additional online resources or studies related to the Arm Cortex-A76
  \item gem5 defaults or explicitly documented approximations when no public implementation-specific value was available
\end{enumerate}

Because many commercial microarchitectural parameters are not publicly disclosed,
several values required approximation. These assumptions are documented in
Table~\ref{tab:model_parameters}. The table is intended to make the model
reproducible and to distinguish between parameters derived from public
documentation and parameters selected by approximation.

\begin{table}[htbp]
\centering
\footnotesize
\setlength{\tabcolsep}{3pt}
\renewcommand{\arraystretch}{1.08}
\caption{Summary of gem5 model parameters and sources.}
\label{tab:model_parameters}
\begin{tabularx}{\textwidth}{ll l l X}
\toprule
\textbf{Cat.} & \textbf{Parameter} & \textbf{Value} & \textbf{Source} & \textbf{Notes} \\
\midrule
CPU & Cores & 4 & \cite{RPIDocs} & Raspberry Pi 5 core count \\
CPU & Frequency & 2400 MHz & \cite{RPIDocs} & Fixed to hardware maximum \\
CPU & Fetch width & 4 & \cite{ArmCommunity} &  \\
CPU & Decode width & 4 & Approx. & Matched to fetch width \\
CPU & Rename width & 4 & Approx. & Matched to fetch width \\
CPU & Dispatch width & 8 & \cite{ArmCommunity} &  \\
CPU & Issue width & 8 & Approx. & Matched to dispatch width \\
CPU & Writeback width & 8 & Approx. & Matched to dispatch width \\
CPU & Commit width & 8 & Approx. & Matched to dispatch width \\
CPU & ROB entries & 128 & \cite{Android_Authority_2018} &  \\
CPU & Load queue & 68 & \cite{Android_Authority_2018} &  \\
CPU & Store queue & 72 & \cite{Android_Authority_2018} &  \\
CPU & Interrupts & Default & \cite{armCortexA76Product} & A76 uses GIC \\
\midrule
L1 & Size & 64 KiB & \cite{armCortexA76Product} & Per-cache value \\
L1 & Associativity & 4 & \cite{armCortexA76Product} &  \\
L1 & Latency & 4 cycles & \cite{Android_Authority_2018} &  \\
L1 & Replacement & TreePLRURP & \cite{armCortexA76Product} & Pseudo-LRU approximation \\
L1D & Prefetcher & StridePrefetcher & \cite{FetchBench,armCortexA76Product} & Approximate L1D prefetching \\
\midrule
L2 & Size & 512 KiB & \cite{armCortexA76Product} &  \\
L2 & Latency & 9 cycles & \cite{Android_Authority_2018} &  \\
L2 & MSHRs & 46 & \cite{A76G76} &  \\
L2 & Replacement & WeightedLRURP & \cite{armCortexA76Product} & Approx. for dynamic biased RP \\
L2 & Prefetcher & StridePrefetcher & \cite{FetchBench,armCortexA76Product} & Approximate L2 prefetching \\
\midrule
L3 & Size & 2 MiB & \cite{Raspberry_Pi_5_2026} & BCM2712 system cache \\
L3 & Latency & 31 cycles & \cite{Android_Authority_2018} & Approximate \\
L3 & MSHRs & 94 & \cite{A76G76} & Approximate \\
\midrule
ITLB & Entries & 48 & \cite{armCortexA76Product} & Fully associative \\
DTLB & Entries & 48 & \cite{armCortexA76Product} & Fully associative \\
L2 TLB & Entries & 1280 & \cite{armCortexA76Product} & 5-way associative \\
\midrule
Memory & Capacity & 8 GiB & \cite{Raspberry_Pi_5_2026} & Same as hardware \\
Memory & DRAM model & DDR4-2400 4x16 & \cite{Raspberry_Pi_5_2026} & Closest gem5 option to LPDDR4X-4267 \\
\bottomrule
\end{tabularx}
\end{table}



\section{Benchmark Suite}

Three benchmarks were selected: breadth-first search (BFS), general matrix
multiplication (GEMM), and the fast Fourier transform (FFT). These workloads were
chosen to represent different computational and memory access patterns that are
important for embedded and edge devices. The justification for each algorithm
comes from the categories discussed in \textit{The Landscape of Parallel
Computing Research: A View from Berkeley}\cite{berkley-view}. That paper
identifies graph traversal, dense linear algebra, and spectral methods as
important computational patterns.

\subsection{Breadth-First Search}

BFS represents an irregular graph traversal workload. Its computational
complexity is commonly expressed as \(O(V + E)\), where \(V\) is the number of
vertices and \(E\) is the number of edges. BFS is typically more memory-bound
than compute-bound, and its performance depends heavily on the regularity of the
graph structure. This study uses the GraphBLAS library, an open-source graph
library built on linear algebra operations\cite{GraphBLAS}. In this case, the
BFS subroutine is implemented using sparse matrix operations.

\subsection{General Matrix Multiplication}

GEMM is a dense linear algebra operation with computational complexity \(O(N^3)\)
for square matrices. GEMM has a regular memory access pattern and high arithmetic
intensity when implemented efficiently. This makes GEMM the most
compute-bound algorithm of the three tested. This study uses OpenBLAS, an
optimized Basic Linear Algebra Subprograms library.

\subsection{Fast Fourier Transform}

The FFT is a spectral method. For an input of size \(N\), its computational
complexity is \(O(N \log N)\). FFT has a more structured access pattern than BFS
but is less straightforward than GEMM. Depending on the implementation, the FFT
may involve strided memory access, butterfly operations, bit-reversal
permutations, and varying locality across stages. This study uses FFTW, a highly
tuned CPU implementation of the FFT.

\section{Experimental Design}

Each benchmark was executed under a sweep of problem sizes. GEMM and FFT were
also tested with varying thread counts. BFS was unable to complete execution when
run in parallel, so only single-threaded BFS results are reported. Problem size
was varied to observe how model accuracy changes with computation size, memory
access behavior, and cache pressure. Thread count was varied to evaluate
multicore scaling behavior and cache coherence accuracy.

For each workload, the same benchmark source was compiled for execution on both
the Raspberry Pi 5 and gem5. All binaries were compiled statically, with compiler
optimization and vectorization enabled where relevant. All parallel code was
compiled using OpenMP.

\subsection{Measurement Collection}

The primary comparison metric was benchmark runtime. Runtime was selected because
it is a general indicator of performance and is commonly used in prior simulation
validation studies\cite{Butko2012-vg,Endo2014-ca}.

Additional architectural statistics were collected to explain performance
differences and determine where the gem5 simulation does not yet match the true
Cortex-A76 implementation. These statistics include IPC, branch miss rate, L1
data miss rate, and L2 data miss rate. In gem5, statistics were collected
directly from the simulator output. On physical hardware, comparable measurements
were collected using Linux \texttt{perf}.

\subsection{Changes from Original Proposal}

The final implementation followed the original proposal closely, with two main
deviations. First, the final model used gem5 \texttt{system-call emulation}
rather than \texttt{full-system} mode. This affected the final compiled binaries.
OpenBLAS made the \textit{mbind} system call, which is unimplemented in
\texttt{system-call emulation} mode. OpenBLAS also attempted to access the
\textit{cntvct\_el0} and \textit{cntfrq\_el0} registers, which are not available
in \texttt{system-call emulation} mode. Second, the first two proposed
milestones---producing a baseline Arm CPU and configuring an out-of-order Arm
CPU---were combined into a single CPU iteration.

Therefore, the final methodology consists of three major stages:

\begin{enumerate}
  \item Construct a runnable Arm gem5 model using the \texttt{O3} CPU model.
  \item Tune CPU, cache, memory, TLB, and related parameters using public documentation and justified approximations.
  \item Evaluate the model against Raspberry Pi 5 hardware using BFS, GEMM, and FFT across varying problem sizes and thread counts.
\end{enumerate}

\section{Results}

This section presents runtime results and architectural statistics for all three
benchmarks. Iteration 1 represents the baseline \texttt{ArmO3CPU} model with few
modifications. Iteration 2 represents the tuned CPU implementation described in
Table~\ref{tab:model_parameters}.




\begin{table}[htbp]
\centering
\footnotesize
\setlength{\tabcolsep}{3pt}
\renewcommand{\arraystretch}{1.04}
\caption{Runtime Comparison of Hardware vs Iteration 1 vs Iteration 2.}
\label{tab:runtime_table}
\begin{tabular}{llcccccccccccc}
\toprule
Algorithm & Problem size & \multicolumn{3}{c}{t=1} & \multicolumn{3}{c}{t=2} & \multicolumn{3}{c}{t=3} & \multicolumn{3}{c}{t=4} \\
\cmidrule(lr){3-5}\cmidrule(lr){6-8}\cmidrule(lr){9-11}\cmidrule(lr){12-14}
 & & hardware & gem5 i1 & gem5 i2 & hardware & gem5 i1 & gem5 i2 & hardware & gem5 i1 & gem5 i2 & hardware & gem5 i1 & gem5 i2 \\
\midrule
\multirow{6}{*}{BFS} & 512 & 0.00472 & \makecell{0.001244\\{\scriptsize -73.6\%}} & \makecell{0.002031\\{\scriptsize -57.0\%}} & X & X & X & X & X & X & X & X & X \\
 & 1024 & 0.00542 & \makecell{0.001818\\{\scriptsize -66.5\%}} & \makecell{0.003179\\{\scriptsize -41.3\%}} & X & X & X & X & X & X & X & X & X \\
 & 2048 & 0.00701 & \makecell{0.003139\\{\scriptsize -55.2\%}} & \makecell{0.005809\\{\scriptsize -17.1\%}} & X & X & X & X & X & X & X & X & X \\
 & 4096 & 0.00934 & \makecell{0.005039\\{\scriptsize -46.0\%}} & \makecell{0.009617\\{\scriptsize +3.0\%}} & X & X & X & X & X & X & X & X & X \\
 & 8192 & 0.01528 & \makecell{0.01039\\{\scriptsize -32.0\%}} & \makecell{0.01928\\{\scriptsize +26.2\%}} & X & X & X & X & X & X & X & X & X \\
 & 16384 & 0.02648 & \makecell{0.02438\\{\scriptsize -7.9\%}} & \makecell{0.04062\\{\scriptsize +53.4\%}} & X & X & X & X & X & X & X & X & X \\
\midrule
\multirow{7}{*}{FFT} & 128 & 0.00035 & \makecell{0.000419\\{\scriptsize +19.7\%}} & \makecell{0.000645\\{\scriptsize +84.3\%}} & 0.00039 & \makecell{0.000496\\{\scriptsize +27.2\%}} & \makecell{0.000733\\{\scriptsize +87.9\%}} & 0.00043 & \makecell{0.000512\\{\scriptsize +19.1\%}} & \makecell{0.000765\\{\scriptsize +77.9\%}} & 0.00043 & \makecell{0.000508\\{\scriptsize +18.1\%}} & \makecell{0.000754\\{\scriptsize +75.3\%}} \\
 & 256 & 0.00028 & \makecell{0.000429\\{\scriptsize +53.2\%}} & \makecell{0.000657\\{\scriptsize +134.6\%}} & 0.00038 & \makecell{0.000503\\{\scriptsize +32.4\%}} & \makecell{0.00074\\{\scriptsize +94.7\%}} & 0.00042 & \makecell{0.000518\\{\scriptsize +23.3\%}} & \makecell{0.000772\\{\scriptsize +83.8\%}} & 0.00042 & \makecell{0.000513\\{\scriptsize +22.1\%}} & \makecell{0.000759\\{\scriptsize +80.7\%}} \\
 & 512 & 0.00031 & \makecell{0.000449\\{\scriptsize +44.8\%}} & \makecell{0.000687\\{\scriptsize +121.6\%}} & 0.00045 & \makecell{0.000516\\{\scriptsize +14.7\%}} & \makecell{0.000756\\{\scriptsize +68.0\%}} & 0.00044 & \makecell{0.000529\\{\scriptsize +20.2\%}} & \makecell{0.000785\\{\scriptsize +78.4\%}} & 0.00044 & \makecell{0.000523\\{\scriptsize +18.9\%}} & \makecell{0.00077\\{\scriptsize +75.0\%}} \\
 & 1024 & 0.0004 & \makecell{0.000493\\{\scriptsize +23.2\%}} & \makecell{0.000747\\{\scriptsize +86.8\%}} & 0.00045 & \makecell{0.000532\\{\scriptsize +18.2\%}} & \makecell{0.00077\\{\scriptsize +71.1\%}} & 0.00048 & \makecell{0.000533\\{\scriptsize +11.0\%}} & \makecell{0.000776\\{\scriptsize +61.7\%}} & 0.00046 & \makecell{0.000527\\{\scriptsize +14.6\%}} & \makecell{0.000766\\{\scriptsize +66.5\%}} \\
 & 2048 & 0.00049 & \makecell{0.000604\\{\scriptsize +23.3\%}} & \makecell{0.000912\\{\scriptsize +86.1\%}} & 0.00057 & \makecell{0.000596\\{\scriptsize +4.6\%}} & \makecell{0.000861\\{\scriptsize +51.1\%}} & 0.00051 & \makecell{0.000581\\{\scriptsize +13.9\%}} & \makecell{0.000843\\{\scriptsize +65.3\%}} & 0.00051 & \makecell{0.000567\\{\scriptsize +11.2\%}} & \makecell{0.000819\\{\scriptsize +60.6\%}} \\
 & 4096 & 0.00077 & \makecell{0.000904\\{\scriptsize +17.4\%}} & \makecell{0.001338\\{\scriptsize +73.8\%}} & 0.00066 & \makecell{0.000769\\{\scriptsize +16.5\%}} & \makecell{0.001089\\{\scriptsize +65.0\%}} & 0.00065 & \makecell{0.000711\\{\scriptsize +9.4\%}} & \makecell{0.001009\\{\scriptsize +55.2\%}} & 0.00065 & \makecell{0.000671\\{\scriptsize +3.2\%}} & \makecell{0.00095\\{\scriptsize +46.2\%}} \\
 & 8192 & 0.0013 & \makecell{0.001313\\{\scriptsize +1.0\%}} & \makecell{0.001825\\{\scriptsize +40.4\%}} & 0.00098 & \makecell{0.001028\\{\scriptsize +4.9\%}} & \makecell{0.001409\\{\scriptsize +43.8\%}} & 0.00087 & \makecell{0.000924\\{\scriptsize +6.2\%}} & \makecell{0.00129\\{\scriptsize +48.3\%}} & 0.00107 & \makecell{0.000855\\{\scriptsize -20.1\%}} & \makecell{0.001207\\{\scriptsize +12.8\%}} \\
\midrule
\multirow{5}{*}{GEMM} & 64 & 0.0003 & \makecell{0.000253\\{\scriptsize -15.7\%}} & \makecell{0.000344\\{\scriptsize +14.7\%}} & 0.00029 & \makecell{0.000253\\{\scriptsize -12.8\%}} & \makecell{0.000344\\{\scriptsize +18.6\%}} & 0.00029 & \makecell{0.000253\\{\scriptsize -12.8\%}} & \makecell{0.000344\\{\scriptsize +18.6\%}} & 0.00029 & \makecell{0.000253\\{\scriptsize -12.8\%}} & \makecell{0.000344\\{\scriptsize +18.6\%}} \\
 & 128 & 0.00149 & \makecell{0.001277\\{\scriptsize -14.3\%}} & \makecell{0.001925\\{\scriptsize +29.2\%}} & 0.00089 & \makecell{0.000696\\{\scriptsize -21.8\%}} & \makecell{0.001024\\{\scriptsize +15.1\%}} & 0.00058 & \makecell{0.000638\\{\scriptsize +10.0\%}} & \makecell{0.00096\\{\scriptsize +65.5\%}} & 0.0005 & \makecell{0.000415\\{\scriptsize -17.0\%}} & \makecell{0.000582\\{\scriptsize +16.4\%}} \\
 & 256 & 0.00984 & \makecell{0.00916\\{\scriptsize -6.9\%}} & \makecell{0.01401\\{\scriptsize +42.4\%}} & 0.00568 & \makecell{0.004674\\{\scriptsize -17.7\%}} & \makecell{0.007119\\{\scriptsize +25.3\%}} & 0.00405 & \makecell{0.003605\\{\scriptsize -11.0\%}} & \makecell{0.005381\\{\scriptsize +32.9\%}} & 0.00279 & \makecell{0.00249\\{\scriptsize -10.8\%}} & \makecell{0.00372\\{\scriptsize +33.3\%}} \\
 & 512 & 0.07948 & \makecell{0.0804\\{\scriptsize +1.2\%}} & \makecell{0.1144\\{\scriptsize +44.0\%}} & 0.04091 & \makecell{0.04354\\{\scriptsize +6.4\%}} & \makecell{0.05863\\{\scriptsize +43.3\%}} & 0.0291 & \makecell{0.02938\\{\scriptsize +1.0\%}} & \makecell{0.04279\\{\scriptsize +47.0\%}} & 0.06561 & \makecell{0.02809\\{\scriptsize -57.2\%}} & \makecell{0.0342\\{\scriptsize -47.9\%}} \\
 & 1024 & 0.6215 & \makecell{0.6083\\{\scriptsize -2.1\%}} & \makecell{0.8864\\{\scriptsize +42.6\%}} & 0.3215 & \makecell{0.3382\\{\scriptsize +5.2\%}} & \makecell{0.4504\\{\scriptsize +40.1\%}} & 0.2385 & \makecell{0.2204\\{\scriptsize -7.6\%}} & \makecell{0.3116\\{\scriptsize +30.7\%}} & 0.5127 & \makecell{0.238\\{\scriptsize -53.6\%}} & \makecell{0.2427\\{\scriptsize -52.7\%}} \\
\bottomrule
\end{tabular}
\end{table}

The tuned model was ultimately unsuccessful in producing a consistently more
accurate recreation of the Arm Cortex-A76 CPU architecture. Across benchmarks,
Iteration 2, which contained the tuned parameters, was less similar to physical
hardware than the baseline Iteration 1, with the exception of single-threaded
BFS. Across all tests, problem sizes, and thread counts, Iteration 1 differed
from the physical Arm Cortex-A76 by an average of 16.8\%. In contrast, Iteration
2 differed from physical hardware by an average of 52.6\%.

It is also important to note that the difference between hardware and simulation
varied more by algorithm than by thread count. Across all problem sizes, thread
counts, and implementations, FFT had an average difference of 44.8\%, while GEMM
had an average difference of 24.6\%. By comparison, when grouped by thread count
for FFT and GEMM, thread counts 1 through 4 had average differences of 40.0\%,
31.7\%, 32.3\%, and 34.8\%, respectively. The fact that error varied more by
algorithm than by thread count suggests that workload-specific memory behavior
was a larger source of mismatch than multicore scaling alone. The stronger GEMM
results also support this interpretation, since GEMM has a more regular and
cache-friendly access pattern than FFT or BFS.

That said, there was no strong correlation between runtime difference and L1 data
miss rate. In fact, no strong trends appeared in L1 miss rate across algorithms,
problem sizes, or thread counts. This may be due to differences between how
\texttt{perf} and gem5 define and collect cache miss statistics. As a result, the
L1 data miss rate measurements are useful as supporting evidence but should not
be treated as directly interchangeable between hardware and simulation.

One important observation from development is that, for much of the tuning
process, the second CPU iteration performed closer to hardware than the baseline
iteration. Near the end of experimentation, cache latency values were introduced
that produced the first results in which Iteration 2 was consistently farther
from hardware. This suggests that the model is sensitive to cache and memory
latency assumptions and that the final parameter set should not be interpreted as
a definitive tuned Cortex-A76 model.

Another limiting factor is the lack of public information about the Arm
Cortex-A76 branch predictor. Because implementation-specific branch predictor
information was unavailable, the predictor was left at its gem5 default. As a
consequence, the average difference between the hardware branch miss rate and the
simulated branch miss rate was 618\%. This large mismatch indicates that branch
prediction is likely another major source of model error.

\section{Conclusion}
This study set out to evaluate whether gem5 could be used to construct a useful
approximation of the Raspberry Pi 5's Broadcom BCM2712 CPU subsystem. A baseline
Arm \texttt{O3} model and a second tuned model were constructed and compared
against physical Raspberry Pi 5 hardware using BFS, GEMM, and FFT. The results
show that the baseline model often tracked hardware runtime more closely than the
tuned model, while the tuned model improved some single-threaded BFS results but
introduced larger errors across FFT and GEMM.

The main conclusion is that matching individual public Cortex-A76 parameters does
not automatically produce a more accurate full-system performance model. The
final tuned model was highly sensitive to cache latency, memory hierarchy, and
branch prediction assumptions. These components are also among the least publicly
documented parts of the Cortex-A76 implementation. As a result, the validation
results suggest that gem5 can provide a useful experimental framework for
studying broad performance trends, but that a Cortex-A76 model should be
validated iteratively against hardware rather than assumed correct after parameter
matching.

Future work should focus on isolating the source of the mismatch more carefully.
A useful next step would be to tune one subsystem at a time, beginning with cache
and memory latency, then evaluating branch prediction and prefetching separately.
Additional microbenchmarks designed to stress one architectural feature at a time
would make it easier to distinguish between CPU pipeline error, cache hierarchy
error, DRAM modeling error, and library or system-call-emulation artifacts.
Moving selected experiments to gem5 \texttt{full-system} mode may also reduce
some of the limitations introduced by \texttt{system-call emulation}. Overall,
this work establishes a reproducible starting point for BCM2712 modeling and
shows that validation against physical hardware is essential when using gem5 for
architecture studies.

\bibliographystyle{ACM-Reference-Format}
\bibliography{sample-base}

\end{document}

